\documentclass[11pt]{article}
\usepackage[english]{babel}
\newcommand{\imgdir}{images/en}
% Shared layout for the Illuminati manual. Each manual-<lang>.tex loads this
% after \documentclass and sets \manuallang (babel name) and \imgdir.
% Build with XeLaTeX (see build.ps1); the fonts are the app's own.

\usepackage{fontspec}
\setmainfont{Literata}[
  Path=../../assets/fonts/literata/, Extension=.ttf,
  UprightFont=*-Regular, BoldFont=*-Bold, ItalicFont=*-Italic, BoldItalicFont=*-BoldItalic]
\setsansfont{AtkinsonHyperlegible}[
  Path=../../assets/fonts/atkinson/, Extension=.ttf,
  UprightFont=*-Regular, BoldFont=*-Bold, ItalicFont=*-Italic, BoldItalicFont=*-BoldItalic]
\newfontfamily\dyslexicfont{OpenDyslexic}[
  Path=../../assets/fonts/opendyslexic/, Extension=.otf,
  UprightFont=*-Regular, BoldFont=*-Bold, ItalicFont=*-Italic, BoldItalicFont=*-Bold-Italic]
\newfontfamily\garamondfont{EBGaramond}[
  Path=../../assets/fonts/ebgaramond/, Extension=.ttf,
  UprightFont=*-Regular, BoldFont=*-Bold, ItalicFont=*-Italic, BoldItalicFont=*-BoldItalic]

\usepackage[a4paper, margin=2.4cm, headheight=14pt]{geometry}
\usepackage{graphicx}
\usepackage[table]{xcolor}
\usepackage{booktabs, tabularx, array}
\usepackage{enumitem}
\usepackage{fancyhdr}
\usepackage{titlesec}
\usepackage[most]{tcolorbox}
\usepackage[hidelinks]{hyperref}

% Colours from the app's logo.
\definecolor{violet}{HTML}{34206A}
\definecolor{violetlight}{HTML}{EDE8F7}
\definecolor{gold}{HTML}{C9A227}
\definecolor{goldlight}{HTML}{FBF4DF}

\setlength{\parskip}{0.55em}
\setlength{\parindent}{0pt}
\setlist{itemsep=0.25em, topsep=0.3em}
\linespread{1.08}

\titleformat{\section}{\sffamily\Large\bfseries\color{violet}}{\thesection}{0.8em}{}
\titleformat{\subsection}{\sffamily\large\bfseries\color{violet}}{\thesubsection}{0.7em}{}
\titlespacing*{\section}{0pt}{1.6em}{0.6em}
\titlespacing*{\subsection}{0pt}{1.1em}{0.35em}
% Don't leave a heading alone at the bottom of a page.
\usepackage{needspace, etoolbox}
\pretocmd{\section}{\needspace{8\baselineskip}}{}{}
\pretocmd{\subsection}{\needspace{5\baselineskip}}{}{}

\pagestyle{fancy}
\fancyhf{}
\fancyhead[L]{\sffamily\small\color{violet} Illuminati}
\fancyhead[R]{\sffamily\small\color{gray}\nouppercase{\rightmark}}
\fancyfoot[C]{\sffamily\small\thepage}
\renewcommand{\headrulewidth}{0.3pt}
% \rightmark is the first section mark on the page.
\renewcommand{\sectionmark}[1]{\markboth{#1}{#1}}
\renewcommand{\subsectionmark}[1]{}

% Keep catalogue names whole in every language.
\babelhyphenation{Wiki-source Guten-berg}

% A key or key combination, e.g. \kbd{Ctrl+F}.
\newcommand{\kbd}[1]{\tikz[baseline=(k.base)]\node[draw=gray!60, fill=gray!8,
  rounded corners=2pt, inner sep=2pt, font=\sffamily\small](k){#1};}
\usepackage{tikz}

% Arrow keys (the sans font has no arrow characters).
\newcommand{\kright}{\kbd{\ensuremath{\rightarrow}}}
\newcommand{\kleft}{\kbd{\ensuremath{\leftarrow}}}
\newcommand{\kup}{\kbd{\ensuremath{\uparrow}}}
\newcommand{\kdown}{\kbd{\ensuremath{\downarrow}}}

% A button, menu item or label as it appears in the app.
\newcommand{\ui}[1]{\textsf{\textbf{#1}}}

% Screenshot with caption; skipped if the image file isn't there.
\newcommand{\shot}[2]{%
  \IfFileExists{\imgdir/#1.jpg}{%
    \begin{figure}[htbp]\centering
      \fcolorbox{gray!40}{white}{\includegraphics[width=0.96\textwidth]{\imgdir/#1.jpg}}
      \caption{#2}
    \end{figure}}{}}

\newtcolorbox{tip}[1]{enhanced, colback=goldlight, colframe=gold,
  boxrule=0.6pt, arc=3pt, left=8pt, right=8pt, top=5pt, bottom=5pt,
  fonttitle=\sffamily\bfseries, coltitle=black, colbacktitle=gold!35,
  title={#1}}
\newtcolorbox{note}[1]{enhanced, colback=violetlight, colframe=violet,
  boxrule=0.6pt, arc=3pt, left=8pt, right=8pt, top=5pt, bottom=5pt,
  fonttitle=\sffamily\bfseries, coltitle=white, colbacktitle=violet,
  title={#1}}

% Two-column tables that fill the text width.
\newenvironment{keytable}[2]{%
  \renewcommand{\arraystretch}{1.3}%
  \rowcolors{2}{white}{violetlight}%
  \tabularx{\textwidth}{@{}>{\raggedright\arraybackslash}p{0.33\textwidth}X@{}}
  \toprule \textsf{\textbf{#1}} & \textsf{\textbf{#2}} \\ \midrule}{%
  \bottomrule\endtabularx}

% Title page: \maketitlepage{subtitle}{version line}{blurb}
\newcommand{\maketitlepage}[3]{%
  \begin{titlepage}
    \centering
    \vspace*{3cm}
    \includegraphics[width=4.2cm]{../../assets/icon/app_icon.png}\par\vspace{1.2cm}
    {\sffamily\fontsize{40}{46}\selectfont\bfseries\color{violet} Illuminati\par}
    \vspace{0.5cm}
    {\sffamily\LARGE #1\par}
    \vspace{1.5cm}
    {\large\itshape #3\par}
    \vfill
    {\sffamily\color{gray} #2\par}
  \end{titlepage}}

\hypersetup{pdftitle={Illuminati — User manual}, pdfauthor={Illuminati}, pdflang={en}}

\begin{document}

\maketitlepage{User manual}{Version 1.0 · Windows, macOS, Linux and Android}%
  {Find, download and read books that were never under copyright\\or whose copyright has expired.}

{\setcounter{tocdepth}{1}\tableofcontents}
\newpage

% ---------------------------------------------------------------------------
\section{Welcome}

Illuminati is a reading app for Windows, macOS, Linux and Android. It works like the e-book
apps you may know, with one difference: every book it offers is free and legal to download,
because its copyright has expired or it was never under copyright.

\begin{itemize}
  \item Over 75,000 books from \textbf{Project Gutenberg}, plus \textbf{Standard Ebooks},
        \textbf{Wikisource} and the \textbf{Internet Archive}.
  \item A comfortable reader with highlights, notes, bookmarks, search, a dictionary and
        translation.
  \item Your own library, with shelves and reading lists.
  \item Display and accessibility settings, including a dyslexia-friendly font.
  \item No account, no subscription and no advertising. Everything is stored on your device.
\end{itemize}

The interface is available in English, Spanish, French, German, Italian and Portuguese.

% ---------------------------------------------------------------------------
\section{Getting started}

\subsection{Installing}
\begin{itemize}
  \item \textbf{Windows:} copy the whole \texttt{Illuminati} folder to your computer and run
        \texttt{illuminati.exe}. Keep the files next to it; the program needs them.
  \item \textbf{Android:} install the \texttt{.apk} file. You may need to allow installing apps
        from this source in your phone's settings.
  \item \textbf{macOS and Linux:} install the build made for your system.
\end{itemize}

\subsection{The home screen}
On a computer the app opens maximized. The home screen has two parts:

\begin{enumerate}
  \item \ui{Your library} (top): the books you have downloaded, most recently read first, with
        a progress bar. The \ui{Library \& shelves} button opens your whole library.
  \item \ui{Discover free books} (below): a search box, catalog language, categories, a row for
        each extra catalog and, at the bottom, the Project Gutenberg catalog.
\end{enumerate}

\shot{home}{The home screen.}

The two buttons at the top right open \ui{Display \& accessibility} (section~\ref{sec:display})
and \ui{About}, which shows where books come from and the licences of the fonts.

% ---------------------------------------------------------------------------
\section{Finding books}

\subsection{Searching}
Type in \ui{Search by title or author}. The search:
\begin{itemize}
  \item matches any part of a title or author name: \emph{miserab} finds \emph{Les Misérables};
  \item ignores accents and capitals: \emph{misérable} and \emph{MISERABLE} give the same results;
  \item searches every catalog at once; each catalog shows its matches in its own row.
\end{itemize}
Clear the box with the \ui{×} button to return to the popular books.

\subsection{Categories and language}
The chips under the search box (\ui{Fiction}, \ui{Mystery}, \ui{Science Fiction}, \ui{Poetry}…)
narrow every catalog to one subject. The language menu on the right shows only books in one
language, or \ui{All languages}. A catalog with nothing for the chosen language or category is
hidden.

\subsection{The catalogs}
\begin{keytable}{Catalog}{What it offers}
  Project Gutenberg & The largest free library, over 75,000 books, sorted by popularity. Often
                      available in an illustrated edition. \\
  Standard Ebooks   & Carefully proofread classics with beautiful typography, in English. \\
  Wikisource        & Texts transcribed by volunteers, in many languages. The e-book is prepared
                      when you download it; very long works (for example novels in several
                      volumes) may not be available. \\
  Internet Archive  & Scans of old printed editions, published 95 or more years ago. The text was
                      made automatically from the scans and may contain errors. \\
\end{keytable}

\subsection{A book's page}
Tap a cover to open the book's page. It shows the summary, subjects, language and catalog, and a
note about copyright. From here you can:
\begin{itemize}
  \item \ui{Download free}: the book is saved on your device, ready to read offline. A small
        download button also appears on every cover in the catalog.
  \item \ui{Read now} or \ui{Continue reading}, once it is downloaded.
  \item \ui{Want to read}, \ui{Finished} and \ui{Add to shelf}, even before downloading
        (section~\ref{sec:library}).
  \item \ui{Get illustrated edition}, for Gutenberg books you downloaded without pictures.
\end{itemize}

\shot{details}{A book's page.}

% ---------------------------------------------------------------------------
\section{Your library}\label{sec:library}

\ui{Library \& shelves} on the home screen opens every book you have saved: downloaded books and
books you saved for later without downloading (shown with a dimmed cover and a cloud).

\shot{library}{The library, with its views and shelves.}

\subsection{Views, filter and sorting}
The chips at the top show a view, each with the number of books:
\ui{All}, \ui{Downloaded}, \ui{Want to read}, \ui{Reading}, \ui{Finished}, and then each of your
shelves. The box above them filters by title or author. The sort button (top right) orders the
books by \ui{Recently read}, \ui{Title}, \ui{Author} (by surname), \ui{Progress} or
\ui{Recently added}. The app remembers your choice.

\subsection{Reading lists}
\begin{itemize}
  \item \ui{Want to read} saves a book for later, without downloading it. Opening the book removes
        it from this list.
  \item \ui{Finished} is set automatically when you reach the end of a book (you can undo it), or
        by hand.
\end{itemize}

\subsection{Shelves}
Shelves are your own collections, such as \emph{Favorites} or \emph{Philosophy}. A book can be on
several shelves.
\begin{itemize}
  \item Create one with \ui{New shelf} in the library, or from \ui{Add to shelf} on any book.
  \item Right-click (or long-press) a shelf to rename or delete it. Deleting a shelf keeps its
        books in your library.
\end{itemize}

\subsection{The book menu}
Right-click a book (or long-press it on a phone) for \ui{Book details}, \ui{Add to shelf},
\ui{Want to read}, \ui{Mark as finished} and two ways of removing it:
\begin{itemize}
  \item \ui{Remove download} frees the space but keeps the book in your library, with its notes,
        shelves and reading position.
  \item \ui{Remove from library} removes everything, including notes and bookmarks.
\end{itemize}

% ---------------------------------------------------------------------------
\section{Reading}

Tap a downloaded book to open it. It opens where you left off.

\shot{reader}{The reader with its toolbar.}

\subsection{Turning pages}
\begin{keytable}{Action}{How}
  Next page        & Tap the right edge · \kright{} · \kbd{Page Down} · \kbd{Space} \\
  Previous page    & Tap the left edge · \kleft{} · \kbd{Page Up} \\
  Scroll a little  & \kup{} \kdown{} or the mouse wheel \\
  Show the toolbar & Tap the middle of the page \\
\end{keytable}

\subsection{The toolbar}
Tap the middle of the page to show it. At the top: back to the library, search
(section~\ref{sec:search}), bookmark, contents and display settings. At the bottom: previous and
next chapter, a slider to jump between chapters, and your progress.

\subsection{Progress and time left}
The bottom line shows how many minutes are left in the chapter and how far you are in the book.
The time is based on your own reading speed, which the app learns as you read.

\subsection{Contents and pictures}
\ui{Contents} lists the book's real chapters, even when a file holds several of them. Pictures
appear in the text; tap one to see it full screen and zoom in.

% ---------------------------------------------------------------------------
\section{Highlights, notes and bookmarks}

\subsection{Selecting text}
On a phone, long-press a word and drag the handles. On a computer, drag over the text or
double-click a word, then \textbf{right-click} the selection. A menu appears with:
\begin{itemize}
  \item \ui{Highlight}: marks the text in your highlight colour.
  \item \ui{Note}: a highlight with your own note, and a choice of four colours.
  \item \ui{Define} and \ui{Translate} (section~\ref{sec:lookup}).
  \item \ui{Copy}.
\end{itemize}
Highlights snap to whole words. Highlights with a note are underlined. Select a highlighted
passage again to edit its note or remove it.

\subsection{Bookmarks}
Use the bookmark button in the toolbar, or \kbd{Ctrl+B}. A red ribbon shows that the page is
bookmarked.

\subsection{The side panel}
The \ui{Contents} button opens a panel with three tabs: \ui{Contents}, \ui{Bookmarks} and
\ui{Notes}. Tap any entry to jump to it. \ui{Copy all as Markdown} copies all your highlights and
notes, ready to paste into another program.

% ---------------------------------------------------------------------------
\section{Dictionary and translation}\label{sec:lookup}

Select a word (up to three words) and choose \ui{Define} to see its definitions. Definitions in
the book's language are shown first. They come from Wiktionary and are written in English.

Select any passage and choose \ui{Translate}. Choose the target language from the list; the app
remembers it.

\begin{note}{Internet needed}
Definitions and translations need an internet connection. The selected text is sent to
Wiktionary or MyMemory; nothing else is sent.
\end{note}

% ---------------------------------------------------------------------------
\section{Searching inside a book}\label{sec:search}

Press \kbd{Ctrl+F} or use the search button in the toolbar, and type at least two letters. The
search ignores capitals and accents. Choose a result to jump to it. A bar at the top shows
\emph{3 of 27}: use its arrows, or \kbd{F3}, to go to the next result. \kbd{Esc} closes it.

% ---------------------------------------------------------------------------
\section{Display and accessibility}\label{sec:display}

Open it with the sliders button on the home screen, or the text button in the reader. A preview
at the top shows each change as you make it.

\shot{settings}{Display and accessibility.}

\subsection{Language}
\ui{App language} switches the interface to English, Spanish, French, German, Italian or
Portuguese at once. \ui{Same as device} follows your device's language.

\subsection{Text}
\begin{itemize}
  \item \ui{Theme}: light, sepia or dark.
  \item \ui{Font}: four fonts included with the app, so books look the same on every device:
  \begin{itemize}
    \item \textbf{Literata}, designed for reading on screens (the default);
    \item {\garamondfont EB Garamond}, a classic book face;
    \item {\sffamily Atkinson Hyperlegible}, designed so that similar letters are easy to tell
          apart;
    \item {\dyslexicfont OpenDyslexic}, designed for readers with dyslexia.
  \end{itemize}
  \item \ui{Text size}, \ui{Line spacing} and \ui{Letter spacing}.
  \item \ui{Justify text}: turn it off for even spacing between words, which many readers with
        dyslexia find easier.
\end{itemize}

\subsection{Page}
\ui{Margins} and \ui{Line width}. Shorter lines are easier to follow.

\subsection{Screen}
\ui{Brightness} dims the whole app. \ui{Warmth} reduces blue light, like a night mode. Both work
on top of your device's own settings.

\ui{Reset to defaults} restores all these settings.

% ---------------------------------------------------------------------------
\section{Keyboard shortcuts}

\begin{keytable}{Keys}{Action}
  \kright{} \kbd{Page Down} \kbd{Space} & Next page \\
  \kleft{} \kbd{Page Up}                 & Previous page \\
  \kup{} \kdown{}                        & Scroll a little \\
  \kbd{Ctrl+F}                         & Search in the book \\
  \kbd{F3}                             & Next search result \\
  \kbd{Ctrl+B}                         & Bookmark the page, or remove the bookmark \\
  \kbd{Esc}                            & Close search or the toolbar, or go back to the library \\
  \kbd{Shift} + arrows                 & Extend a text selection \\
\end{keytable}

% ---------------------------------------------------------------------------
\section{Privacy and your data}

Illuminati has no account and no tracking. Your library, notes and settings are stored only on
your device (on Windows, in \texttt{\%APPDATA\%\textbackslash Illuminati}).

The app connects to the internet only to:
\begin{itemize}
  \item browse and search the catalogs and download books and covers;
  \item define or translate text you selected, which is sent to Wiktionary or MyMemory.
\end{itemize}

% ---------------------------------------------------------------------------
\section{About copyright}

All books are in the public domain in the United States. Some Wikisource texts are instead under
free licences (such as CC BY-SA), which allow you to read and share them freely. From the Internet
Archive, only books published 95 or more years ago are listed.

\begin{tip}{Check your local laws}
Copyright terms differ between countries. A book that is free in the United States may still be
protected where you live, so check your local laws before downloading.
\end{tip}

% ---------------------------------------------------------------------------
\section{Troubleshooting}

\begin{keytable}{Message}{What to do}
  Couldn't reach a catalog & The service may be busy or your connection down. Use \ui{Try again}
                             later; the other catalogs keep working. \\
  This title has no e-book edition & It is probably an audio book. Look for another edition. \\
  Wikisource could not prepare this book & The work is too long for Wikisource's service. Try
                             one of its volumes, or another catalog. \\
  The free daily translation limit has been reached & MyMemory allows a limited number of free
                             translations per day. Try again tomorrow. \\
  The download took too long & Check your connection and try again. \\
\end{keytable}

% ---------------------------------------------------------------------------
\section{Credits}

Books: Project Gutenberg, Standard Ebooks, Wikisource (via WS Export) and the Internet Archive.
Definitions: Wiktionary (CC BY-SA). Translations: MyMemory.

Fonts, under the SIL Open Font License: Literata, EB Garamond, Atkinson Hyperlegible (Braille
Institute) and OpenDyslexic (Abbie Gonzalez). Their licences are under \ui{About} → \ui{View
licenses} in the app.

\end{document}
